\section{Appendix}

\subsection{Framework Comparison}
\label{app:framework_comparison}

\begin{table}[H]
\centering
\scriptsize
\setlength{\tabcolsep}{4.5pt}
\renewcommand{\arraystretch}{1.15}
\begin{tabular}{lcccc}
\toprule
\textbf{System} &
\makecell{\textbf{Async RL}\\\textbf{Support}} &
\makecell{\textbf{Async}\\\textbf{Rollout Staging}} &
\makecell{\textbf{Rollout as}\\\textbf{Service}} &
\makecell{\textbf{Agent Harness}\\\textbf{Agnostic}} \\
\midrule
\polar & \cmark & \cmark & \cmark & \cmark \\
\prorlagent \citep{prorlagent2026} & \cmark & \cmark & \cmark & \xmark \\
SkyRL-Agent \citep{skyrlagent2025} & \cmark & \cmark & \xmark & \pmark \\
PRIME-RL \citep{primerl2026} & \cmark & \xmark & \xmark & \xmark \\
Agent Lightning \citep{agentlightning2025} & \pmark & \xmark & \pmark & \pmark \\
rLLM \citep{rllm2025} & \pmark & \xmark & \xmark & \xmark \\
OpenClaw-RL \citep{openclawrl2026} & \cmark & \xmark & \xmark & \pmark \\
\bottomrule
\end{tabular}
\caption{\textbf{Comparing rollout-system design choices.} We find that modern rollout infrastructures should meet following criteria: async RL support means training can consume rollouts while generation continues under explicit policy-version or staleness handling; async rollout staging means rollout execution is decomposed into independently scheduled runtime-preparation, execution, post-run reconstruction/evaluation, and cleanup stages; rollout as service means a durable task API that is separable from a specific trainer loop; and native-harness agnosticism means a CLI, SDK, or application harness can be trained without being reimplemented as the framework's environment. \cmark{} denotes first-class support, \pmark{} denotes partial or planned support, and \xmark{} denotes that we did not find the property as a primary design contract in the referenced code or documentation.}
\label{tab:framework_comparison}
\end{table}

The partial marks in \cref{tab:framework_comparison} avoid treating adjacent mechanisms as absent. SkyRL exposes custom generators and Harbor integration, but native harnesses are not the default unit of rollout. Agent Lightning provides a rollout store, queue, runner control plane, and broad framework instrumentation, while its main boundary is trace/workflow observability rather than staged execution of opaque harness processes. rLLM includes fully asynchronous training and a model gateway that captures token IDs and log probabilities, but its rollout service abstraction is narrower than a distributed runtime lifecycle service. OpenClaw-RL decouples serving, rollout collection, judging, and training for real-world agent settings; its support is organized around OpenClaw and specific terminal, GUI, SWE, and tool-call recipes rather than arbitrary native harness submission.

\subsection{SWE-Gym GRPO Hyperparameters}
\label{app:swegym_grpo_hparams}

\begin{table}[H]
\centering
\small
\setlength{\tabcolsep}{6pt}
\renewcommand{\arraystretch}{1.12}
\begin{tabular}{ll}
\toprule
\textbf{Hyperparameter} & \textbf{Value} \\
\midrule
Base checkpoint & \texttt{Qwen/Qwen3.5-4B} \\
Training data & \texttt{NovaSky-AI/SkyRL-v0-293-data}, train split, 293 tasks \\
Trainer & Slime asynchronous GRPO \\
Epochs & 1 \\
Rollout batch size & 4 \\
Samples per prompt & 16 \\
Trace construction & \texttt{prefix\_merging} \\
Optimizer & Adam \\
Learning rate & $1\times10^{-6}$ \\
Weight decay & 0.1 \\
TIS & Enabled \\
\bottomrule
\end{tabular}
\caption{\textbf{Training hyperparameters for the SWE-Gym GRPO experiments.} The table reports ordinary policy-optimization and rollout parameters from \texttt{examples/swegym\_slime\_grpo}; cluster topology and worker placement are omitted.}
\label{tab:swegym_grpo_hparams}
\end{table}

\subsection{Representative Task Payload}
\label{app:task_payload}

\begin{polarlisting}{Representative Polar Task Payload}
{
  "task_id": "polar-swegym-{rollout_id}-{group_index}",
  "instruction": "Fix the issue in /polar/session/workspace.",
  "num_samples": 8,
  "timeout_seconds": 1200,
  "runtime": {
    "backend": "docker",
    "image": "{sample.metadata.runtime_image}",
    "network": "host",
    "workdir": "/polar/session/workspace",
    "prepare": [
      {
        "type": "exec",
        "command": "prepare repository, harness, and dependencies"
      }
    ]
  },
  "agent": {
    "harness": "codex",
    "model_name": "{served_model_name}"
  },
  "builder": {
    "strategy": "prefix_merging"
  },
  "evaluator": {
    "strategy": "swebench_harness",
    "refresh_runtime": true,
    "config": {
      "repo_dir": "/testbed",
      "patch_command": "cd /polar/session/workspace && git add -A && git diff --cached --binary",
      "instance": "{sample.metadata.instance}"
    }
  },
  "callback_url": "http://{trainer_host}:{callback_port}/callback/task_result",
  "metadata": {
    "group_id": "{rollout_group_id}",
    "policy_version": "{policy_version}",
    "rollout_step": "{rollout_step}"
  }
}
\end{polarlisting}

\subsection{Representative Trace}
\label{app:trace_example}

\begin{polarlisting}{Representative Trainer-Facing Trace}
{
  "prompt_ids": [151644, 872, "..."],
  "response_ids": [9211, 374, 264, "...", 151645, 271, 151644],
  "loss_mask": [1, 1, 1, "...", 0, 0, 1],
  "response_logprobs": [
    {
      "token": "The",
      "token_id": 785,
      "logprob": -0.21
    },
    {
      "token": " patch",
      "token_id": 10042,
      "logprob": -0.44
    },
    {
      "token_id": 151645,
      "logprob": 0.0
    }
  ],
  "prompt_messages": [
    {"role": "system", "content": "agent harness system prompt"},
    {"role": "user", "content": "task instruction and current tool state"}
  ],
  "response_messages": [
    {"role": "assistant", "content": "sampled assistant turn"}
  ],
  "tools": null,
  "finish_reason": "stop",
  "reward": 1.0,
  "metadata": {
    "session_id": "session-123",
    "task_id": "polar-swegym-0001",
    "builder": "prefix_merging",
    "harness": "codex"
  }
}
\end{polarlisting}

\subsection{Service API Summary}

The rollout service exposes a small asynchronous API:

\begin{itemize}
    \item \texttt{POST /rollout/task/submit}: submit a non-blocking task request.
    \item \texttt{GET /rollout/task/\{task\_id\}}: poll task status, partial results, and final results.
    \item \texttt{GET /rollout/status}: inspect task states, node states, and pending sessions.
    \item \texttt{POST /callbacks/session\_result}: receive gateway session callbacks.
    \item \texttt{POST /nodes/register} and \texttt{POST /nodes/\{node\_id\}/heartbeat}: maintain gateway membership and scheduling metrics.
\end{itemize}

The gateway exposes a control surface for session creation, status, and deletion, plus a catch-all proxy surface for provider-style model requests. Session deletion is used by the rollout pipeline as best-effort cleanup after a terminal result has been persisted.
